\section{The full ring properties of the original ordinary Hecke ideal}
\label{sec:ring-properties-v44}

\status{This chapter proves the literal \texttt{TJRing} consumer, including
the zero-module branch and all eight fields of \texttt{RingProps}.  The
ordinary Hecke algebra, its scalar map, the universal determinant and the
ideal are the original ones.  The arithmetic inputs are the already
constructed ordinary algebra OHA11, density AFS5, and the actual
\texttt{Out} supplied by HEP12.  We prove the auxiliary-place inertia
assertion needed by this consumer, rather than adding it as a new input.
No finite-rank bound, semisimplicity assumption at that auxiliary place,
or replacement of the coefficient ring is introduced.}

Fix the original prime \(\ell\), coefficient object \(c\), and
\(O=c.O\subset\overline{\mathbb Q}_\ell\).  Thus \(O\) is a DVR
finite over \(\mathbb Z_\ell\).  Write \(k=O/\mathfrak m_O\).
For a frame \(F\), put \(G=F.G_0\), with its given involution
\(g\mapsto g^c\) and cyclotomic character \(\epsilon\).  Let
\(r\) be the original residual datum of rank \(n\),
\(B\) its supplied universal ordinary determinant ring, and
\(\Lambda\) its supplied weight object.  We write
\[
 D=B.u.D,\qquad \psi_{v,i}=B.u.\psi(v,i),\qquad
 \lambda=\Lambda.\mathrm{toPair}(B.u):\Lambda\longrightarrow B.
 \tag{JRP.1}
\]
For an ideal \(J\subset B\), put \(A=B/J\),
\(D_J=D\otimes_B A\), and \(\psi^J_{v,i}=\psi_{v,i}\otimes_B A\).
We use \([g]\) for the basis element of a group algebra.  The original
anti-involution and ordered ordinary element are
\[
 [g]^*=\epsilon(g)^{1-n}[(g^c)^{-1}],\qquad
 o_v(\psi^J;g_0,\ldots,g_{n-1})=
       \prod_{i=0}^{n-1}\bigl([\phi_v(g_i)]-\psi^J_{v,i}(g_i)[1]\bigr),
 \tag{JRP.2}
\]
where the factors are written in increasing \(i\), without commuting
them.  A degree-\(n\) determinant \(D_A\) is the original natural
multiplicative homogeneous polynomial law on the group algebra.
Its kernel predicate has the full meaning
\[
 \begin{gathered}
 D_A.\mathrm{Kills}(y)\quad\Longleftrightarrow\quad
 D_{A,C}(1+h y_C)=1\\
 \hbox{for every commutative \(A\)-algebra \(C\) and every \(h\in C[G]\)}.
 \end{gathered}
 \tag{JRP.3}
\]
Thus equality of traces at group elements alone is not what we descend.

The residual characters are \(\bar\chi_i:G\to k^\times\), and
\[
 I^p_u=I_u\cap\bigcap_i\ker\bar\chi_i.
\]
The original \(\mathrm{V0ok}\) is precisely
\[
 \bar\chi_i|_{I_0}=1\ (\hbox{all }i),\qquad
 \bar\chi_i|_{D_0}=\bar\chi_j|_{D_0}\ (\hbox{all }i,j),\qquad
 \exists g\in D_0:\ \bar\epsilon(g)\ne1.
 \tag{JRP.4}
\]
All algebra homomorphisms below preserve the given \(O\)-structure.

\subsection{Integral algebraic points and the actual auxiliary-place inertia}

% BEGIN JRP1
\paragraph{JRP1. The original coefficient ring and its Teichmuller lifts.}
\label{claim:JRP1}
The given \(O\) is the ring of integers of its fraction field
\(E_0\subset\overline{\mathbb Q}_\ell\), which is finite over
\(\mathbb Q_\ell\).  It is complete, its residue field has order
\(q=\ell^f\), and there is a multiplicative section
\(\tau:k\to O\) of reduction, with \(\tau(1)=1\) and
\(\tau(a)^{q-1}=1\) for \(a\ne0\).  In particular
\[
 t_i=\tau\circ\bar\chi_i:G\longrightarrow O^\times
 \tag{JRP.5}
\]
satisfies the original \(r.\mathrm{IsTeich}(t)\), and
\(\bar\chi_i(g)=1\) implies \(t_i(g)=1\).
\begin{proof}
The given embedding makes \(O\) torsion free over
\(\mathbb Z_\ell\); its finite module is therefore free.  Localizing
gives a finite-dimensional domain over \(\mathbb Q_\ell\), hence
the field \(E_0\).  Every element of \(O\) is integral over
\(\mathbb Z_\ell\).  Conversely an element of \(E_0\) integral
over \(\mathbb Z_\ell\) is integral over \(O\), so belongs to the
integrally closed DVR \(O\).  This identifies \(O\) with the
usual integer ring, equivalently the elements of norm at most one.
Finite freeness over the complete ring \(\mathbb Z_\ell\) gives
\(\ell\)-adic completeness.  Since \(\ell\) is a nonzero
nonunit in the DVR, \(\ell O=\mathfrak m_O^e\) for some
\(e>0\); the two adic filtrations are cofinal.  Reduction is a
finite extension of \(\mathbb F_\ell\).

For \(a\in k^\times\), lift it successively to a root of
\(X^{q-1}-1\) modulo \(\mathfrak m_O^N\).  The derivative
\((q-1)X^{q-2}\) is a unit at every such lift.  Given a lift
modulo \(\mathfrak m_O^N\), the equation for a correction in
\(\mathfrak m_O^N/\mathfrak m_O^{N+1}\) is therefore a linear
equation with a unique solution.  Completeness gives a unique root
\(\tau(a)\) reducing to \(a\).  Products of these roots have the
same property for the product of their residues, so uniqueness
proves multiplicativity.  Put \(\tau(0)=0\).  The root reducing
to one is one.  Composing with \(\bar\chi_i\) gives (JRP.5),
with exponent exactly \(\#k-1\), as required in
\(\mathrm{IsTeich}\).  These constructions use the original
coefficient ring throughout.
\end{proof}
% END JRP1

% BEGIN JRP2
\paragraph{JRP2. Algebraic points preserve the prescribed residue map.}
\label{claim:JRP2}
Let \(R\) be an \(O\)-algebra of the original class
\(\mathrm{IsCO}\), and let
\(x:R\to_O\overline{\mathbb Q}_\ell\) be any algebra homomorphism.
Then
\[
 a\in\mathfrak m_R\ \Longrightarrow\ |x(a)|_\ell<1,
 \qquad |x(b)|_\ell\le1\quad(b\in R).
 \tag{JRP.6}
\]
If \(O'\) is the integer ring of a finite coefficient extension
containing \(O\) and \(x(b)\), then
\[
 \overline{x(b)}=j_{k,k'}\bigl(\mathrm{res}_{O,R}(b)\bigr).
 \tag{JRP.7}
\]
Here \(j_{k,k'}:k\hookrightarrow k'\) is the natural residue-field
inclusion; no continuity of \(x\) is assumed.
\begin{proof}
Suppose \(a\in\mathfrak m_R\) and \(|x(a)|_\ell\ge1\).
Then \(\alpha=x(a)\ne0\), and \(\alpha^{-1}\) is integral
over \(O\): it belongs to the integer ring of the finite extension
\(E_0(\alpha)/E_0\), which is integral over \(O\).
Choose a monic equation
\[
 (\alpha^{-1})^d+c_{d-1}(\alpha^{-1})^{d-1}
                    +\cdots+c_0=0,\qquad c_i\in O.
\]
Multiplication by \(\alpha^d\) says that \(x\) kills
\(1+c_{d-1}a+\cdots+c_0a^d\).  This element lies in
\(1+\mathfrak m_R\), so is a unit in the local ring \(R\).
A unital ring map cannot send a unit to zero.  This contradiction
proves the first assertion.

By the \(\mathrm{res\_surj}\) field of \(\mathrm{IsCO}\),
each \(b\in R\) has \(b-b_0\in\mathfrak m_R\) for some
\(b_0\in O\).  Thus \(|x(b)-b_0|_\ell<1\).
JRP1 and the ultrametric inequality give \(|x(b)|_\ell\le1\),
and reduction in any finite coefficient field containing these
elements gives (JRP.7).  The residue of \(b_0\) is exactly
\(\mathrm{res}_{O,R}(b)\).  This argument proves the needed
residue compatibility directly from the original local ring
structure; no analytic boundedness theorem for the entire point is
needed.
\end{proof}
% END JRP2

% BEGIN JRP3
\paragraph{JRP3. The finite prime-to-residue-characteristic inertia calculation.}
\label{claim:JRP3}
Let \(K\) be a finite extension of \(\mathbb Q_p\), with
absolute group \(G_K\) and inertia \(I_K\).  If
\(c:I_K\to\mu_m\subset\overline K^\times\) has open kernel
and \(p\nmid m\), then
\[
 c(g\tau g^{-1})=g(c(\tau))
       \qquad(g\in G_K,\ \tau\in I_K).
 \tag{JRP.8}
\]
In particular a nontrivial character of \(\ell\)-power order,
\(\ell\ne p\), fixed by conjugation by \(g\), forces
\(\epsilon_\ell(g)\equiv1\pmod\ell\).
\begin{proof}
Inertia fixes \(\mu_m\): roots of unity of order prime to \(p\)
are the unique lifts of their residues.  Distinct such roots have
distinct residues, because the derivative of \(X^m-1\) is a unit
at each root.  We first find \(z\ne0\) with
\(\tau(z)=c(\tau)z\).  On the fixed field of \(\ker c\),
the finite group \(I_K/\ker c\) acts faithfully.  Its distinct
field automorphisms are linearly independent.  One proof takes a
relation with the smallest number of nonzero coefficients,
evaluates it on a product, and subtracts a suitable multiple of the
original relation; two distinct automorphisms can be separated by
one field element, contradicting minimality.  Hence the operator
\(\sum_a c(a)^{-1}a\) is nonzero.  Apply it to an element on
which it is nonzero; reindexing the sum proves the desired eigenlaw.

For fixed \(g\) put \(y=g^{-1}(z)/z\).  Every local Galois
automorphism preserves the absolute value, so \(|y|=1\).
As inertia acts trivially on residue fields,
\(|\tau(y)-y|<1\).  The eigenlaw gives
\[
 \tau(y)-y=
 \frac{g^{-1}\bigl(c(g\tau g^{-1})\bigr)-c(\tau)}{c(\tau)}\,y.
\]
The two terms in the numerator are \(m\)-th roots of unity.
Their difference has norm less than one, so their residues, and
hence the roots themselves, agree.  Applying \(g\) proves
(JRP.8).  For the last assertion choose a nontrivial value of
\(c\) and take a power having order \(\ell\).
The cyclotomic definition is
\(g(\zeta)=\zeta^{\epsilon_\ell(g)}\); fixing an element of
order \(\ell\) gives the asserted congruence.  The same proof
applies to any open subgroup of \(G_{\mathbb Q_p}\), using its
finite fixed field.
\end{proof}
% END JRP3

% BEGIN JRP4
\paragraph{JRP4. Scalar residual characteristic polynomials force auxiliary unramifiedness.}
\label{claim:JRP4}
Let \(p\ne\ell\), let \(K/\mathbb Q_p\) be finite, and let
\(O'\) be a complete DVR with finite residue field \(k'\) of
characteristic \(\ell\).  Suppose
\(\rho_0:G_K\to\mathrm{GL}_n(O')\) is continuous modulo every
power of \(\mathfrak m_{O'}\), and
\[
 \det(X-\bar\rho_0(g))=(X-\kappa(g))^n
                  \quad(g\in G_K),
 \tag{JRP.9}
\]
where \(\kappa:G_K\to(k')^\times\) is trivial on inertia.
If some \(g_*\in G_K\) has
\(\bar\epsilon_\ell(g_*)\ne1\), then
\(\rho_0|_{I_K}=1\).  No semisimplicity or bound on \(n\)
is required.
\begin{proof}
Let \(M=\rho_0(g)\), and lift \(\kappa(g)\) to \(b\in O'\).
Cayley--Hamilton applied to (JRP.9) gives
\((\bar M-\bar b I)^n=0\).  Since \(\ell^n\ge n\),
the characteristic-\(\ell\) binomial identity gives
\(M^{\ell^n}\equiv b^{\ell^n}I\pmod{\mathfrak m_{O'}}\).
For \(j\ge1\), if \(M\equiv bI\pmod{\mathfrak m^j}\),
then
\(M^\ell\equiv b^\ell I\pmod{\mathfrak m^{j+1}}\).
Indeed the linear binomial term has a factor \(\ell\in\mathfrak m\),
and every higher term has at least two factors in \(\mathfrak m^j\).
All matrices in this calculation commute with the scalar matrix.
Iteration therefore gives, for each \(N\ge1\),
\[
 \rho_0(g^{\ell^{n+N}})
       \equiv b^{\ell^{n+N}}I\pmod{\mathfrak m_{O'}^N}.
 \tag{JRP.10}
\]

Reduce \(\rho_0\) modulo \(\mathfrak m_{O'}^N\); the target
is finite.  Its inertia image \(H\) has exponent dividing
\(\ell^{n+N}\), by (JRP.10) with \(b=1\).
Cauchy's theorem makes \(H\) an \(\ell\)-group.  If \(H\)
were nontrivial it would have a nontrivial abelian character into
\(\ell\)-power roots of unity.  For completeness, a finite
\(\ell\)-group has nontrivial center by the class equation;
induction on its order gives a nontrivial homomorphism to a cyclic
group of order \(\ell\), unless the group is already abelian,
where the assertion is immediate.  Composing such a character with
the inertia quotient gives a nontrivial
\(c:I_K\to\mu_{\ell^h}\) with open kernel.

The image of \(g_*^{\ell^{n+N}}\) is scalar by (JRP.10), so
it centralizes \(H\) and fixes \(c\).  JRP3 forces
\(\bar\epsilon_\ell(g_*^{\ell^{n+N}})=1\).
But this equals
\(\bar\epsilon_\ell(g_*)^{\ell^{n+N}}\), which cannot be one:
Frobenius is injective on the finite field \(k'\).
Thus \(H=1\).  This holds for every \(N\); separation of the
adic filtration now gives \(\rho_0(\tau)=1\) for every
\(\tau\in I_K\).
\end{proof}
% END JRP4

% BEGIN JRP5
\paragraph{JRP5. The literal auxiliary-place assertion for every original point.}
\label{claim:JRP5}
In the original Hida setting let \(n\ge1\), assume
\(r.\mathrm{V0ok}\), and let
\(x:B\to_O\overline{\mathbb Q}_\ell\) and
\(\rho:F.G_0\to\mathrm{GL}_n(\overline{\mathbb Q}_\ell)\)
satisfy
\[
 \rho\text{ is continuous},\qquad
                 D\otimes_{B,x}\overline{\mathbb Q}_\ell=\det\rho.
 \tag{JRP.11}
\]
Then \(\rho(g)=1\) for every \(g\in F.I_0\).
This is the complete original
\texttt{point\_\allowbreak rep\_\allowbreak unramified\_\allowbreak at\_\allowbreak auxiliary\_\allowbreak place} endpoint.
\begin{proof}
We first explain the coefficient issue.  On any compact local
decomposition group, ADK1 puts the coefficients of the continuous
representation in one finite extension of \(\mathbb Q_\ell\).
Enlarge that field to contain \(E_0\), without changing \(O\)
or the point \(x\).  Compactness bounds all matrix entries.
Consequently
\[
 \mathcal L=\sum_g\rho(g)O'^n
\]
lies between \(O'^n\) and \(\varpi^{-M}O'^n\) for some
\(M\).  It is a finitely generated torsion-free module over the
DVR, hence free of rank \(n\).  A basis conjugates the
representation into \(\mathrm{GL}_n(O')\), continuously modulo
each \(\mathfrak m_{O'}^N\).  Conjugation preserves characteristic
polynomials.  JRP2 applied to every coefficient of
\(D.\chi([g])\), together with \(B.u.\mathrm{red}\), proves
that the residual characteristic polynomial of this lattice is
\(\prod_i(X-\bar\chi_i(g))\), transported through \(k\to k'\).
It is the original residue map, even though this argument temporarily
uses a larger field to choose a stable lattice.

Here is the exact local dictionary at the auxiliary place.  Put
\(F^+=L^+\), let \(v_0\) be the prescribed place, and let
\(p\ne\ell\) be its residue characteristic from the Hida setting.
Choose the specified valuation above \(v_0\).  The decomposition
embedding identifies its local absolute group with an open subgroup
of \(G_{\mathbb Q_p}\), preserving inertia and cyclotomic
characters.  The Hida setting puts this decomposition group in the
image of \(G_L\to G_{L^+}\).  Thus its composite with the
quotient map gives a continuous homomorphism
\(\psi_0:G_{L^+_{v_0}}\to F.G_0\) whose image is \(F.D_0\)
and whose inertia image is \(F.I_0\).  More explicitly,
\(\psi_0(t)=V.\mathrm{pr}(e)\) whenever a decomposition element
\(e\in G_L\) restricts to \(t\); independence follows from
the original quotient, and the formula for its underlying element
is the continuous map \(\mathrm{toGS}(t)\).  This is the source
local bridge, not a newly chosen Galois quotient.

Restrict (JRP.11) along \(\psi_0\).  Since \(n\ge1\), choose
the first residual character.  Formula (JRP.4) makes the local
residual polynomial
\[
                  (X-\kappa(t))^n,
       \qquad\kappa=\bar\chi_0\circ\psi_0,
\]
with \(\kappa\) trivial on inertia.  The last clause of (JRP.4)
supplies a local element with nontrivial residual cyclotomic value;
the residue-field inclusion \(k\to k'\) preserves this
inequality.  JRP4 therefore kills the full local inertia.
Its image is precisely \(F.I_0\), proving the claim.  Neither
the global nor the local representation was assumed semisimple.
\end{proof}
% END JRP5

\subsection{Descent of the full determinant laws}

% BEGIN JRP6
\paragraph{JRP6. Jointly separating points detect polynomial laws and their kernels.}
\label{claim:JRP6}
Let \(A\) be a commutative ring and
\(f_s:A\to K_s\) a family of ring maps such that
\(\bigcap_s\ker f_s=0\).  Equality of determinants of the same
degree over \(A\) can be tested after every \(f_s\).
Moreover, if \(D_A\otimes f_s=\det\rho_s\) and
\(\rho_s(y)=0\) for every \(s\), with coefficients of
\(y\in A[G]\) mapped through \(f_s\), then
\(D_A.\mathrm{Kills}(y)\) in the full sense of (JRP.3).
The corresponding coefficientwise test holds for any natural
polynomial identity, including the polarization identity.
\begin{proof}
Coefficientwise, the maps
\(A[X_1,\ldots,X_m]\to K_s[X_1,\ldots,X_m]\) jointly
separate: every coefficient of a polynomial killed by all of them
is zero.  Given a finite list of group elements \(g_1,\ldots,g_m\),
apply the two determinant laws to the generic element
\(\sum_jX_j[g_j]\).  Naturality identifies coefficient
specialization of the resulting polynomials with the corresponding
laws over \(K_s\).  Their equality at all \(s\) therefore gives
equality over \(A[X_1,\ldots,X_m]\).  Specializing \(X_j\) to
arbitrary elements of an arbitrary \(A\)-algebra proves equality
on every group-algebra element, since each has finite support.

For the kernel assertion, use the polynomial
\[
 D_{A,A[\mathbf X]}
       \left(1+\left(\sum_jX_j[g_j]\right)y\right)-1.
 \tag{JRP.12}
\]
After \(f_s\), the matrix of \(y\) is zero, so this polynomial
is zero.  All its coefficients vanish by separation.  Specialize
the variables to the finitely many coefficients of any
\(h\in C[G]\); naturality proves exactly (JRP.3).
For polarization use instead
\(D((\sum_jX_j[g_j])^*)-D(\sum_jX_j[g_j])\).
The scalar cyclotomic coefficients in \(*\) map compatibly, and
the same argument applies.  No injectivity after tensoring with an
arbitrary algebra, no flatness, and no reducedness assumption is
used.  If the index family is empty, its separation hypothesis
forces \(A=0\), and the same coefficient proof remains valid.
\end{proof}
% END JRP6

% BEGIN JRP7
\paragraph{JRP7. The original pairing gives the entire star identity.}
\label{claim:JRP7}
Suppose \(\rho:G\to\mathrm{GL}_n(K)\) and
\(P\in\mathrm{GL}_n(K)\) satisfy the original
\(F.\mathrm{IsForm}(O,n,\rho,P)\):
\[
 \rho(g)^{\mathsf T}P\rho(g^c)=\epsilon(g)^{1-n}P.
 \tag{JRP.13}
\]
Then, for every commutative \(K\)-algebra \(C\) and every
\(y\in C[G]\),
\[
 \rho_C(y)^{\mathsf T}P_C=P_C\rho_C(y^*),\qquad
       \det\rho_C(y^*)=\det\rho_C(y).
 \tag{JRP.14}
\]
In particular \(\rho(y)=0\) implies \(\rho(y^*)=0\).
\begin{proof}
Multiply (JRP.13) on the right by \(\rho(g^c)^{-1}\).
The result is (JRP.14)'s first equality for \([g]\), with
exactly the exponent \(1-n\) in (JRP.2).  Scalar extension
preserves it; linearity gives it for every finite group-algebra
sum.  Taking determinants and canceling the unit \(\det P_C\)
proves the second identity.  When the left matrix is zero,
invertibility of \(P\) proves the kernel assertion.  No symmetry
or alternating property of \(P\) is required by this argument
or by the original \(\mathrm{RingProps}\).
\end{proof}
% END JRP7

% BEGIN JRP8
\paragraph{JRP8. A flag gives both original ordinary relations, in their order.}
\label{claim:JRP8}
Let \(\rho:H\to\mathrm{GL}_n(K)\) have a full flag with
successive characters \(\xi_0,\ldots,\xi_{n-1}\), in the
literal \(\mathrm{HasFlag}\) convention of upper triangular
matrices.  Then
\[
 \det(X-\rho(h))=\prod_{i=0}^{n-1}(X-\xi_i(h)),
 \qquad
 \prod_{i=0}^{n-1}\bigl(\rho(h_i)-\xi_i(h_i)I\bigr)=0
 \tag{JRP.15}
\]
for every \(h\in H\) and every ordered tuple \((h_i)\).
\begin{proof}
Conjugate once to the upper triangular basis supplied by
\(\mathrm{HasFlag}\).  The first equality is the determinant
of an upper triangular matrix.  Let \(F_j\) be the span of
its first \(j\) basis vectors.  Each
\(\rho(h_i)-\xi_i(h_i)I\) preserves the flag and sends
\(F_{i+1}\) into \(F_i\): the action on
\(F_{i+1}/F_i\) is zero.  The rightmost factor \(i=n-1\)
therefore sends the entire space into \(F_{n-1}\); the next
factor sends this into \(F_{n-2}\), and so on until \(F_0=0\).
This is precisely the ascending written order in (JRP.15).
Conjugating back proves both statements.  For rank zero the
vector space and its identity endomorphism are zero, so the empty
matrix product has the same assertion.  No factors with different
indices were commuted.
\end{proof}
% END JRP8

% BEGIN JRP9
\paragraph{JRP9. The quotient by the specified ideal has separating points.}
\label{claim:JRP9}
Suppose maps \(x_s:B\to_O K\) satisfy
\[
 J\subset\ker x_s\quad(\hbox{all }s),\qquad
       \bigl[x_s(b)=0\ (\hbox{all }s)\bigr]\Longrightarrow b\in J.
 \tag{JRP.16}
\]
They induce jointly separating maps \(f_s:A=B/J\to_O K\).
If \(D\otimes x_s=\det\rho_s\), then
\(D_J\otimes f_s=\det\rho_s\), and the specialized ordinary
characters are \(f_s\psi^J_{v,i}=x_s\psi_{v,i}\).
\begin{proof}
Define \(f_s(\bar b)=x_s(b)\).  The first clause of
(JRP.16) proves independence of the representative.  If every
\(f_s(\bar b)\) is zero, the second clause gives \(b\in J\),
hence \(\bar b=0\).  The quotient and specialization maps
compose to \(x_s\).  Naturality of determinant scalar extension
and the entrywise definition of pushing a unit-valued character
give the final two equalities.  These are equalities on the actual
quotient \(B/J\), rather than on a reduced or characteristic-zero
replacement of it.
\end{proof}
% END JRP9

% BEGIN JRP10
\paragraph{JRP10. Descent of the full polarization field.}
\label{claim:JRP10}
In JRP9, assume each \(\rho_s\) has a matrix satisfying
\(F.\mathrm{IsForm}(O,n,\rho_s,P_s)\).  Then
\(F.\mathrm{IsPolar}(O,D_J)\), with its universal scalar-algebra
quantifier.  Furthermore any \(y\in A[G]\) whose matrix is zero
at all these points satisfies both
\(D_J.\mathrm{Kills}(y)\) and \(D_J.\mathrm{Kills}(y^*)\).
\begin{proof}
JRP7 proves the polarization polynomial identity after each
\(f_s\).  JRP6 descends it coefficientwise to \(A\), then to
all its scalar algebras.  This is exactly \(\mathrm{IsPolar}\).
For the second assertion, JRP7 makes the specialized matrices of
both \(y\) and \(y^*\) zero.  Apply JRP6's kernel test to each.
The proof retains the universal definition of \(\mathrm{Kills}\)
and does not infer it merely from a characteristic polynomial at
\(y\).
\end{proof}
% END JRP10

% BEGIN JRP11
\paragraph{JRP11. Descent of O1 and O2 for the same ordinary characters.}
\label{claim:JRP11}
In JRP9--JRP10, suppose every local restriction
\(\rho_s\circ\phi_v\) has \(\mathrm{HasFlag}\) with the
ordered characters \(x_s\psi_{v,i}\).  Then
\[
 D_J.\chi([\phi_v(g)])=
                  \prod_i\bigl(X-\psi^J_{v,i}(g)\bigr),
 \tag{JRP.17}
\]
and for every original tuple \(g_i\in F.G_v\), both
\(o_v(\psi^J;\mathbf g)\) and
\(o_v(\psi^J;\mathbf g)^*\) are killed by \(D_J\).
\begin{proof}
After each \(f_s\), the characteristic polynomial of the
determinant is the characteristic polynomial of the actual matrix
\(\rho_s(\phi_v(g))\).  JRP8 makes it the indicated product.
Every coefficient of the difference in (JRP.17) is killed by all
\(f_s\), so is zero by JRP9.  This is O1, with the original
pushforward characters.

Specializing (JRP.2) gives the ordered product of matrices in the
second formula of (JRP.15), because coefficient maps preserve each
individual factor and its order.  This product is zero.  JRP10
therefore kills both the original ordinary element and its star.
This is O2 in full, for arbitrary tuples, not only the tuple whose
entries are all equal.
\end{proof}
% END JRP11

% BEGIN JRP12
\paragraph{JRP12. Descent of every prescribed type, O3, and the guarded auxiliary field.}
\label{claim:JRP12}
In JRP9--JRP10, suppose that for every original Teichmuller
system \(t\) satisfying \(r.\mathrm{IsTeich}(t)\), every
\(s,u\), one matrix simultaneously diagonalizes
\(\rho_s|_{I_u}\) as \(\operatorname{diag}(t_i|_{I_u})\).
Then the full \(\mathrm{types}\) and O3 fields of
\(\mathrm{RingProps}\) hold.  If also
\(r.\mathrm{V0ok}\Rightarrow\rho_s|_{I_0}=1\) for every
\(s\), its \(\mathrm{unr0}\) field holds.
\begin{proof}
For any supplied \(t\), conjugation does not change the
determinant polynomial law on \(K[I_u]\).  Thus every
\(f_s\) identifies the restriction of \(D_J\) with the
determinant of the diagonal representation whose entries are the
images of \(t_i\) under \(O\to A\).  JRP6 descends equality
of the entire laws.  This proves \(\mathrm{types}\) for every
\(t\), retaining its quantifier before \(u\).

For O3 use the specific system constructed in JRP1.  If
\(g\in I^p_u\), all \(\bar\chi_i(g)=1\), so all
\(t_i(g)=1\).  The supplied diagonalization then gives
\(\rho_s(g)=1\) for every \(s\).  Apply JRP10 to
\(y=[g]-1\).  Since \(*\) is additive and fixes one,
\(y^*=[g]^*-1\); these are exactly the two O3 kernel
statements.  Under the additional displayed guard, the same
argument for \(g\in I_0\) proves \(\mathrm{unr0}\).
The guard \(\mathrm{V0ok}\) has not been dropped.
\end{proof}
% END JRP12

% BEGIN JRP13
\paragraph{JRP13. Transfer through the original surjection, with the identical kernel.}
\label{claim:JRP13}
Let \(f:B\twoheadrightarrow T\) be the original
\(O\)-algebra map, with \(J=\ker f\), and suppose
\[
                   f(\lambda(a))=\mathrm{scal}_T(a)
                                  \quad(a\in\Lambda).
 \tag{JRP.18}
\]
Assume \(T\) is finite and torsion free over \(\Lambda\),
and a family of good maps \(z:T\to_O K\) jointly separates
\(T\).  If their composite points \(x=z\circ f\) have the
representations, forms, flags, types and guarded inertia properties
of JRP10--JRP12, then \(B.\mathrm{RingProps}(\Lambda,J)\).
\begin{proof}
For this family of composite points, (JRP.16) follows directly:
all points kill \(\ker f\), and vanishing at them forces
\(f(b)=0\) by separation in \(T\).  JRP10--JRP12 therefore
give \(\mathrm{polar}\), O1, O2, O3,
\(\mathrm{types}\) and \(\mathrm{unr0}\) on \(B/J\).

The surjection gives an actual \(O\)-algebra isomorphism
\(e:B/J\simeq T\), \(e(\bar b)=f(b)\).
Equation (JRP.18) makes it \(\Lambda\)-linear for the
original scalar map \(a\mapsto\overline{\lambda(a)}\).
Pulling back a finite generating set for \(T\) proves exactly
the \(\mathrm{fin}\) field.  If \(a\ne0\) and
\(\lambda(a)b\in J\), apply \(f\) to get
\(\mathrm{scal}_T(a)f(b)=0\).  Torsion freeness gives
\(f(b)=0\), hence \(b\in J\), which is precisely
the \(\mathrm{tf}\) field.  No new ideal was defined in this
argument, and no passage to a fraction field was used to establish
these two integral conclusions.
\end{proof}
% END JRP13

% BEGIN JRP14
\paragraph{JRP14. The entire top-ideal branch.}
\label{claim:JRP14}
\leavevmode\par
For the original supplied \(B,\Lambda\), one has
\(B.\mathrm{RingProps}(\Lambda,\top)\).
\begin{proof}
The quotient \(B/\top\) is the zero ring.  Its scalar map from
\(\Lambda\) is finite, generated as a module by the empty set
(or by zero).  The \(\mathrm{tf}\) implication has conclusion
\(b\in\top\), so always holds.  Every commutative algebra over
the zero ring is itself zero, because its unital scalar map sends
\(1=0\) to \(1=0\).  Consequently all values of every
polynomial law agree, every polynomial equality holds, and every
\(\mathrm{Kills}\) assertion is true with its full scalar-algebra
quantifier.  This proves polarization, O1, O2, O3, every type
equality and the guarded \(\mathrm{unr0}\) assertion.  In
particular no nonempty family of arithmetic points is assumed in
this branch.
\end{proof}
% END JRP14

% BEGIN JRP15
\paragraph{JRP15. The full original TJRing endpoint.}
\label{claim:JRP15}
\leavevmode\par
Assume exactly the original \(H'=\mathrm{HypothesesH'}\),
the Hida data \(\ell,L,\Omega,s,n\) with \(n\ge1\),
\(V,Q,c,r,\iota,B,\Lambda\), an original
\(R:\mathrm{UnitaryHeckeData}(Q,\iota,\Lambda)\), and a
level \(K\), with \(V.S'=s.S\).  For every ideal
\(J\subset B\),
\[
             R.\mathrm{IsJ}(B,K,J)
                    \quad\Longrightarrow\quad
             B.\mathrm{RingProps}(\Lambda,J).
 \tag{JRP.19}
\]
\begin{proof}
Expand the original \(\mathrm{IsJ}\) disjunction.  If the
ordinary residual module is zero, its clause gives \(J=\top\),
so JRP14 applies.  Otherwise retain its actual nonzero module,
the given commutativity and \(\mathrm{IsCO}\) structures, and
its surjection
\[
 f:B\twoheadrightarrow
 T=(\mathrm{hset}(Q,R.\tau_0)).TK(R.\theta,R.a,K),
 \qquad J=\ker f,
 \tag{JRP.20}
\]
together with the original \(\mathrm{IsMap}\) witness and
(JRP.18).  No surjection or ideal is reselected.

AFS5 makes this given \(\Lambda\) a domain and supplies
\(\mathrm{DenseWt}(R.\theta,\lambda_0)\) for the zero
dominant weight \(\lambda_0\).  That weight is admitted by
the original \(\mathrm{hset}\).  OHA11 applied at this same
level proves that this same \(T\) is finite over \(\Lambda\)
and that its nonzero scalars act injectively.  The hypotheses
\(\mathrm{IsWt}\), \(\mathrm{IsRes}\) and the actual
ordinary-module axioms are supplied by the fixed \(R\) and
the installation used in HEP12.

Use HEP12's \(\mathrm{Out}\) for the same \(D\otimes f\),
characters \(f\psi_{v,i}\), algebra \(T\), and zero weight.
Its \(\mathrm{sep}\) field, supplied by AFS20, says exactly
that the maps \(z:T\to_O\overline{\mathbb Q}_\ell\) whose
restrictions are good at \(\lambda_0\) jointly separate \(T\).
For every such \(z\), its \(\mathrm{pt}\) field supplies
one actual continuous representation \(\rho_z\), with
\[
 (D\otimes f)\otimes z=\det\rho_z.
\]
Its \(\mathrm{UnitaryPointProperties}\) give the original
form matrix and the simultaneous inertia diagonalizations for
\emph{every} Teichmuller system.  The same point field gives
\(\mathrm{HasFlag}\) with the ordered characters
\(z(f\psi_{v,i})\) under this exact good-point guard.
These are the data used in JRP10--JRP12.  If
\(r.\mathrm{V0ok}\), apply JRP5 to the composite point
\(x=z\circ f\) and to this very \(\rho_z\).  Functoriality
of determinant scalar extension gives (JRP.11), so the missing
auxiliary inertia condition is proved for each selected point.

JRP13 now supplies all eight fields at the original ideal in
(JRP.20).  Their exact dictionary is
\begin{center}
\begin{tabular}{ll}
\texttt{fin}, \texttt{tf} & OHA11, AFS5 and JRP13, over the original \(\Lambda\);\\
\texttt{polar} & the given form and JRP10;\\
\texttt{O1}, \texttt{O2} & the original ordered flags and JRP11;\\
\texttt{O3}, \texttt{types} & every original inertia type and JRP12;\\
\texttt{unr0} & the retained \(\mathrm{V0ok}\) guard, JRP5 and JRP12.
\end{tabular}
\end{center}
The arguments retain \(n\ge1\), every original level and
coefficient object, and the equality \(V.S'=s.S\).
This proves the whole source
\texttt{ring\_props\_of\_is\_j\_unitary\_hecke\_data},
not only its reduction or a property at individual points.
\end{proof}
% END JRP15

\paragraph{Scope.}
JRP15 reconstructs the full \texttt{TJRing} producer with the named
global and ordinary inputs retained.  It proves a component of the
original \texttt{HidaProps}; it does not claim the other arithmetic
purity, classicality, twisting or global counting fields merely from
this ring argument.  The auxiliary assertion JRP5 is independently
available for every original algebraic point satisfying (JRP.11).
